\documentclass{ctexart}

% 数学相关宏包
\usepackage{amsmath,mathtools,amsthm,amssymb}
\usepackage{bbm}
\usepackage{mathrsfs}
\usepackage{centernot}

% 图形和绘图
\usepackage{graphicx,float}
\usepackage{tikz,tikz-cd}
\usepackage{pgfplots}
\usepackage[framemethod=TikZ]{mdframed}
\usetikzlibrary{decorations.pathmorphing,positioning,arrows,automata,calc,shapes,patterns}
\pgfplotsset{compat=1.18}
\usepgfplotslibrary{statistics}

% 表格和列表
\usepackage{booktabs,multirow,colortbl}
\usepackage{enumitem}
\setlist[itemize,enumerate,description]{noitemsep}
\setlistdepth{9}

% 算法
\usepackage{algorithm,algorithmic}

% 字体和颜色
\usepackage{xcolor,listings}

% 页面和链接
\usepackage{fancyhdr,caption}
\usepackage{hyperref}
\usepackage{cleveref}

% 工具宏包
\usepackage{etoolbox,letltxmacro,docmute}
\usepackage{gensymb}

% 自定义设置
\renewcommand{\arraystretch}{1.5}
\let\originalmathbb\mathbb
\renewcommand{\mathbb}[1]{%
    \ifstrequal{#1}{1}{\mathbbm{#1}}{\originalmathbb{#1}}%
}

% 颜色定义
\definecolor{lightblue}{RGB}{173,216,230}
\definecolor{lightgreen}{RGB}{144,238,144}
\definecolor{lightyellow}{RGB}{255,255,224}
\definecolor{lightgray}{RGB}{211,211,211}

% 超链接设置
\hypersetup{
    colorlinks=true,
    linkcolor=red,
    filecolor=magenta,
    urlcolor=cyan,
    pdftitle={数学笔记},
    pdfauthor={乐绎华},
    pdfsubject={数学笔记},
    pdfkeywords={数学},
    bookmarksnumbered=true,
    bookmarksopen=true,
    bookmarksopenlevel=6,
    pdfstartview=FitH,
    pdfpagemode=UseOutlines,
    pdfborder={0 0 0}
}

% 定理环境
\theoremstyle{plain}
\newtheorem{theorem}{定理}[section]
\newtheorem{lemma}[theorem]{引理}
\newtheorem{corollary}[theorem]{推论}
\newtheorem{proposition}[theorem]{命题}

\theoremstyle{definition}
\newtheorem{definition}[theorem]{定义}
\newtheorem{example}[theorem]{例}
\newtheorem{exercise}[theorem]{练习}
\newtheorem{problem}[theorem]{题目}

\theoremstyle{remark}
\newtheorem{remark}[theorem]{备注}
\newtheorem{note}[theorem]{注}
\newtheorem{observation}[theorem]{观察}

% 解答环境
\newtheorem*{solution}{解}
\newtheorem*{proof*}{证明}

% 图片路径
\graphicspath{{F:/iCloudDrive/iCloud~md~obsidian/2025-summer/figures/}}

\author{乐绎华}
\title{数学笔记}
\date{\today}

\begin{document}

\maketitle
\tableofcontents

位势方程 $\Delta u=f$ 在零 Dirichlet 边值条件下的解是泛函 $F(u)=\int_{\Omega}^{} \left[ \frac{1}{2}\lvert \nabla u \rvert ^{2}+fu \right] \, \mathrm{d}x$ 的最小值点，叫做 Dirichlet 原理.

给定方程
\[
-\sum_{i, j=1}^n a_{i j} \partial_{i j}^2 u+\sum_{i=1}^n b_i \partial_i u+c u=f
\]
若 $(a_{ij}(x))$ 总是正定或负定，则称其为椭圆型. 若存在 $c_0>0$，使得
\begin{figure}[H]
\centering
\includegraphics[width=\textwidth]{elliptic-2025072510.png}
% \caption{}
\label{}
\end{figure}
则称为一致椭圆型.

边值条件有：

\begin{itemize}
	\item Dirichlet 边值：$\left.u\right|_{\partial \Omega}(x)=h(x)$.
	\item Neumann 边值：$\left.\frac{ \partial u(x) }{ \partial \mu }\right|_{\partial \Omega}=h(x)$. 其中 $\mu$ 是某个向量场.
	\item Robin 边值：$\frac{ \partial u (x) }{ \partial \mu }+\sigma (x) u (x)=h (x),\forall x\in \partial \Omega$.
	\item ...
\end{itemize}

\begin{figure}[H]
\centering
\includegraphics[width=\textwidth]{elliptic-2025072511.png}
% \caption{}
\label{}
\end{figure}
\begin{figure}[H]
\centering
\includegraphics[width=\textwidth]{1-elliptic-2025072511.png}
% \caption{}
\label{}
\end{figure}
\begin{figure}[H]
\centering
\includegraphics[width=\textwidth]{2-elliptic-2025072511.png}
% \caption{}
\label{}
\end{figure}
\begin{figure}[H]
\centering
\includegraphics[width=\textwidth]{3-elliptic-2025072511.png}
% \caption{}
\label{}
\end{figure}
\begin{proof}
应用紧算子的 Fredholm 定理.
\end{proof}
考虑如下方程
\[
\begin{cases}
Lu=f & \forall x\in \Omega \\
u=\psi & \forall x\in \partial \Omega
\end{cases}
\]
其中 $f\in H^{-1}(\Omega)$, $\psi\in H^{\frac{1}{2}}(\partial\Omega)$.

接下来研究解的正则性问题：如果 $f$ 属于比 $H^{-1}(\Omega)$ 正则性更好的空间，$\psi$ 属于比 $H^{\frac{1}{2}}(\partial \Omega)$ 正则性更好的空间，那么解 $u$ 是否存在更好的正则性？

\begin{figure}[H]
\centering
\includegraphics[width=\textwidth]{4-elliptic-2025072511.png}
% \caption{}
\label{}
\end{figure}
\begin{figure}[H]
\centering
\includegraphics[width=\textwidth]{5-elliptic-2025072511.png}
% \caption{}
\label{}
\end{figure}
略过正则性这些各种范数估计.

接下来考虑特征值问题：
\[
\begin{cases}
Lu=\lambda u & \forall x\in \Omega  \\
u=0 & \forall x\in \partial \Omega 
\end{cases}
\]
我们要研究对于什么样的复数 $\lambda$ ，这个问题存在非零解 $u$，称为\textbf{特征函数}. 所以本节研究的问题叫做 Dirichlet 特征值问题，应用前两节得到的结果和关于紧算子的 Riesz-Shauder 理论，可以得到这个问题的圆满解答.

\begin{figure}[H]
\centering
\includegraphics[width=\textwidth]{elliptic-2025072512.png}
% \caption{}
\label{}
\end{figure}
\begin{figure}[H]
\centering
\includegraphics[width=\textwidth]{1-elliptic-2025072512.png}
% \caption{}
\label{}
\end{figure}
\begin{figure}[H]
\centering
\includegraphics[width=\textwidth]{2-elliptic-2025072512.png}
% \caption{}
\label{}
\end{figure}
这个定理在抽象空间中就有一套证明，参见于品.

接下来考虑极值原理

\begin{figure}[H]
\centering
\includegraphics[width=\textwidth]{3-elliptic-2025072512.png}
% \caption{}
\label{}
\end{figure}
\begin{figure}[H]
\centering
\includegraphics[width=\textwidth]{4-elliptic-2025072512.png}
% \caption{}
\label{}
\end{figure}
\begin{figure}[H]
\centering
\includegraphics[width=\textwidth]{5-elliptic-2025072512.png}
% \caption{}
\label{}
\end{figure}
我们注意到在经典解的极值问题中，证明利用到了经典解的假设，却不涉及解的导数，自然地我们会问：对强解和弱解是否也成立类似的结果？这个问题的回答是肯定的.

\begin{figure}[H]
\centering
\includegraphics[width=\textwidth]{6-elliptic-2025072512.png}
% \caption{}
\label{}
\end{figure}
\begin{figure}[H]
\centering
\includegraphics[width=\textwidth]{7-elliptic-2025072512.png}
% \caption{}
\label{}
\end{figure}
\begin{figure}[H]
\centering
\includegraphics[width=\textwidth]{8-elliptic-2025072512.png}
% \caption{}
\label{}
\end{figure}
\begin{figure}[H]
\centering
\includegraphics[width=\textwidth]{9-elliptic-2025072512.png}
% \caption{}
\label{}
\end{figure}

接下来建立 Dirichlet 边值问题解的 $L^{p}$ 估计，这里先不考虑解的存在性，而是对已知的解建立一些不等式，所以叫做\textbf{先验估计}.

\begin{figure}[H]
\centering
\includegraphics[width=\textwidth]{elliptic-2025072517.png}
% \caption{}
\label{}
\end{figure}
\begin{figure}[H]
\centering
\includegraphics[width=\textwidth]{1-elliptic-2025072521.png}
% \caption{}
\label{}
\end{figure}
\begin{figure}[H]
\centering
\includegraphics[width=\textwidth]{elliptic-2025072521.png}
% \caption{}
\label{}
\end{figure}

\begin{table}[h]
	\centering
	\begin{tabular}{|c|c|}
		\hline
		1 & 1 \\
		\hline
		2 & 2 \\
		\hline
	\end{tabular}
\end{table}

\end{document}